\documentclass[12pt]{article}
%---DOCUMENT MARGINS---
\usepackage{geometry} % Required for adjusting page dimensions and margins
\geometry{
	paper=a4paper, % Paper size, change to letterpaper for US letter size
	top=4cm, % Top margin
	bottom=2cm, % Bottom margin
	left=2cm, % Left margin
	right=2cm, % Right margin
	headheight=3cm, % Header height
	%footskip=1.5cm, % Space from the bottom margin to the baseline of the footer
	headsep=0.5cm, % Space from the top margin to the baseline of the header
	%showframe, % Uncomment to show how the type block is set on the page
}
\usepackage{cmap}

\usepackage[T1,T2A]{fontenc}
\usepackage{fontspec}
\setmainfont[
	BoldFont={* Bold},
	ItalicFont={* Italic},
	BoldItalicFont={* BoldItalic}
]{DejaVu Serif Condensed}
\setsansfont{DejaVu Sans}
\setmonofont{Dejavu Sans Mono}
\usepackage[math-style=TeX]{unicode-math}
\setmathfont{Latin Modern Math}

\usepackage[nobottomtitles*]{titlesec}
\titleformat
{\section} % command
{\fontsize{14}{14}\sffamily\bfseries} % format
{} % label
{0pt} % sep
{} % before-code
\titlespacing{\section}{0pt}{0.75em}{0.25em}
\newcommand{\tl}{}
\newcommand{\ml}{}
\newcommand{\problem}[1]{%
	\section[#1]{#1 \fontsize{14}{14}\selectfont{\hfill{\normalfont{
				\emoji{hourglass-not-done}\tl \space\space \emoji{floppy-disk} \ml}}}}
}
\AddToHook{cmd/section/before}{\clearpage}

\usepackage[dvipsnames]{xcolor}
\titleformat
{\subsection} % command
{\fontsize{14}{14}\sffamily\bfseries} % format
{\includesvg[width=0.035\textwidth, inkscapelatex=false]{lion}\space} % label
{0pt} % sep
{} % before-code
\titlespacing{\subsection}{0pt}{0.75em}{0.25em}

\setlength{\parskip}{0.5em}
\setlength{\parindent}{0pt}
\renewcommand{\baselinestretch}{1.2}
\sloppy

\usepackage{listings}
\definecolor{lstbg}{RGB}{250,250,252}
\definecolor{lstframe}{RGB}{220,220,225}
\lstdefinestyle{mystyle}{
	backgroundcolor=\color{lstbg},
	frame=single,
	rulecolor=\color{lstframe},
	basicstyle=\ttfamily\small,
	keywordstyle=\bfseries,
	breaklines=true,
	aboveskip=0.5\baselineskip,
	belowskip=0\baselineskip  
}
\lstset{style=mystyle, language=C++}

\usepackage{fancyhdr}
\pagestyle{fancy}
\setlength{\headwidth}{\textwidth}
\newcommand{\round}{}
\newcommand{\problemName}{}
\newcommand{\lang}{English}
\fancyhead[L]{
	\begin{minipage}{0.4\textwidth}
		\bf{
			EJOI 2025 \round\\
			\problemName\\
			\emoji{globe-with-meridians} \lang
		}
	\end{minipage}
	\vspace{0.5cm}
}
\fancyhead[R]{%
	\begin{minipage}{0.6\textwidth}
		\includesvg[width=\textwidth, inkscapelatex=false]{logo}
	\end{minipage}
	\vspace{0.5cm}
}
\usepackage{lastpage}
\fancyfoot[C]{\problemName\space(\lang) \thepage\ / \pageref*{LastPage}}

\raggedbottom

\usepackage{amsmath}
\usepackage{stmaryrd}

\usepackage{graphicx}
\graphicspath{{./}}
\usepackage[export]{adjustbox}
\usepackage{wrapfig}
\makeatletter
\patchcmd\WF@putfigmaybe{\lower\intextsep}{}{}{\fail}
\AddToHook{env/wrapfigure/begin}{\setlength{\intextsep}{0pt}}
\makeatother
\usepackage[inkscapearea=page, inkscapepath=./svg-inkscape]{svg}
\svgpath{{./}}

\usepackage{makecell}
\usepackage{tabularray}
\AtBeginEnvironment{table}{\vspace{-0.2cm}}
\AtEndEnvironment{table}{\vspace{-0.2cm}}
\usepackage{float}
\usepackage{placeins}
\usepackage{caption}
\captionsetup[table]{
	skip=0.25em,
	singlelinecheck=false, justification=justified
}

\usepackage{enumitem}
\setlist{itemsep=-0.4em, leftmargin=22pt, topsep=-\parskip}
\AfterEndEnvironment{itemize}{\vspace{\parskip}}
\newcommand{\tabitem}{\indent~~\llap{\textbullet}~~}

\usepackage{hyperref}
\hypersetup{
	colorlinks=true,
	citecolor=blue,
	linkcolor=blue,
	urlcolor=cyan,
}

\usepackage{emoji}

\usepackage[english]{babel}

\begin{document}

\renewcommand{\round}{Day 1}
\renewcommand{\problemName}{ამოცანა ციხე}
\renewcommand{\lang}{ქართული}

\renewcommand{\tl}{$2$ sec.}
\renewcommand{\ml}{$1024$ MB}
\problem{\problemName}

ალისა და ბობი უსამართლოდ გამოამწყვდიეს მაქსიმალურად დაცულ ციხეში. ახლა მათ უნდა შეადგინონ გაქცევის გეგმა. ამის მისაღწევად მათ უნდა მოახერხონ ერთმანეთთან რაც შეიძლება ეფექტური კომუნიკაციის დამყარება (კონკრეტულად, ალისამ უნდა გაუგზავნოს ბობს ინფორმაცია ყოველდღიურად). მათ არ შეუძლიათ ერთმანეთთან შეხვედრა და ინფორმაციას ცვლიან მხოლოდ ხელსახოცებზე დაწერილი მესიჯების გამოყენებით. ალისას უნდა, რომ ბობს ყოველ დღე გაუგზავნოს ახალი ინფორმაცია -- რიცხვი $0$-დან $N-1$-მდე. ყოველ სადილზე ალისა ღებულობს 3 ახალ ხელსახოცს და წერს თითოეულ მათგანზე $0$-დან $M-1$-მდე რიცხვს (არა აუცილებლად განსხვავებულ რიცხვებს) და ტოვებს ხელსახოცებს თავის სკამზე. შემდეგ მათი მოწინააღმდეგე, ჩარლი, ანადგურებს ხელსახოცებიდან ერთ-ერთს და ერთმანეთში ურევს დარჩენილ ორს. საბოლოოდ, ბობი პოულობს სკამზე 2 ხელსახოცს და კითხულობს მათზე რიცხვებს. მან ზუსტად უნდა გაარკვიოს იმ რიცხვის მნიშვნელობა, რომლის გაგზავნაც ალისას სურდა. ხელსახოცებზე სივრცე შეზღუდულია, ამიტომ  $M$ არის ფიქსირებული.  თუმცა, ალისას და ბობის მიზანია ინფორმაციის გამტარუნარიანობის მაქსიმიზაცია, რომ მათ შეძლონ რაც შეიძლება დიდი $N$-ის გაგზავნა. დაეხმარეთ ალისას და ბობს თითოეული მათგანისთვის სტრატეგიის იმპლემენტაციაში, რათა მაქსიმალურად გაზარდონ $N$-ის მნიშვნელობა.

\subsection{იმპლემენტაციის დეტალები}
რადგან ეს კომუნიკაციის ამოცანაა, თქვენი პროგრამა შესრულდება ორ ცალკეულ გამოძახებაში (ერთი ალისასთვის და მეორე ბობისთვის), რომლებსაც არ შეუძლიათ მონაცემების გაზიარება ან კომუნიკაცია ქვემოთ აღწერილის გარდა სხვა გზით. თქვენ იმპლემენტაცია უნდა გაუკეთოთ სამ ფუნქციას:

\begin{lstlisting}
int setup(int M);
\end{lstlisting}

ეს ფუნქცია გამოიძახება ერთხელ ალისას სტრატეგიის შესრულებამდე და ერთხელ ბობის სტრატეგიის შესრულებამდე. მას ეძლევა $M$ და უნდა დააბრუნოს სასურველი $N$, რომლის გადაცემასაც შეეცდება ალისა ბობისთვის. \texttt{setup}-ის ორივე გამოძახებამ უნდა დააბრუნოს ერთი და იგივე $N$.

\begin{lstlisting}
std::vector<int> encode(int A);
\end{lstlisting}

ეს ფუნქცია იმპლემენტაციას უკეთებს ალისას სტრატეგიას. ის გამოიძახება, ბობისთვის $A$ ($0 \leq A < N$) რიცხვის გადასაგზავნად და უნდა დააბრუნოს სამი რიცხვი $W_1, W_2, W_3$ ($0 \leq W_i < M$), რომლებშიც კოდირებულია $A$ რიცხვი. ეს ფუნქცია გამოიძახება სულ $T$-ჯერ -- დღეში ერთხელ ($A$-ს მნიშვნელობები შეიძლება განმეორდეს დღეებს შორის).

\begin{lstlisting}
 int decode(int X, int Y);
\end{lstlisting}

ეს ფუნქცია იმპლემენტაციას უკეთებს ბობის სტრატეგიას. გამოძახებისას ამ ფუნქციას გადმოეცემა \texttt{encode}-ს მიერ დაბრუნებული სამი რიცხვიდან ორი, რაღაც თანმიმდევრობით. მან უნდა დააბრუნოს იგივე მნიშვნელობა $A$, რაც \texttt{encode}-მა მიიღო. ეს ფუნქცია ასევე გამოიძახება $T$-ჯერ --  \texttt{encode}-ის $T$ ცალი გამოძახების შესაბამისად და იმავე თანმიმდევრობით. \texttt{encode}-ის ყველა გამოძახება მოხდება \texttt{decode}-ის ყველა გამოძახებამდე.

\subsection{შეზღუდვები}

\begin{itemize}
\item $M \leq 4300$
\item $T = 5000$
\end{itemize}

\subsection{შეფასება}

კონკრეტული ქვეამოცანისთვის, მიღებული ქულების S წილი დამოკიდებულია \texttt{setup}-ის მიერ ამ ქვეამოცანაში ყველა ტესტს შორის დაბრუნებულ უმცირეს $N$-ზე. ეს ასევე დამოკიდებულია $N^*$-ზე, რომელიც წარმოადგენს $N$-ის სამიზნე მნიშვნელობას, რომელიც გჭირდებათ ქვეამოცანისთვის სრული ქულების მისაღებად:

\begin{itemize}
\item თუ თქვენი ამოხსნა ვერ შეასრულებს რომელიმე ტესტს, მაშინ $S = 0$.
\item თუ $N \geq N^*$, მაშინ $S = 1.0$.
\item თუ $N < N^*$, მაშინ $S = \max\left(0.35 \max\left(\frac{\log(N) - 0.985 \log(M)}{\log(N^*) - 0.985 \log(M)}, 0.0\right)^{0.3} + 0.65 \left(\frac{N}{N^*}\right)^{2.4}, 0.01\right)$.
\end{itemize}

\subsection{ქვეამოცანები}

\begin{table}[H]
\begin{tblr}{|Q[c,m]|Q[c,m]|Q[c,m]|Q[c,m]|}
\hline
\textbf{ქვეამოცანა} & \textbf{ქულა} & $M$ & $N^*$ \\
\hline
$1$ & $10$ & $700$ & $82017$ \\
\hline
$2$ & $10$ & $1100$ & $202217$ \\
\hline
$3$ & $10$ & $1500$ & $375751$ \\
\hline
$4$ & $10$ & $1900$ & $602617$ \\
\hline
$5$ & $10$ & $2300$ & $882817$ \\
\hline
$6$ & $10$ & $2700$ & $1216351$ \\
\hline
$7$ & $10$ & $3100$ & $1603217$ \\
\hline
$8$ & $10$ & $3500$ & $2043417$ \\
\hline
$9$ & $10$ & $3900$ & $2536951$ \\
\hline
$10$ & $10$ & $4300$ & $3083817$ \\
\hline
\end{tblr}
\end{table}
\FloatBarrier

\subsection{მაგალითი}

განვიხილოთ შემდეგი მაგალითი, სადაც $T = 5$. აქ გვაქვს კოდირების სქემა, სადაც ალისა $0$-ის კოდირებისთვის სამ თანაბარ რიცხვს აგზავნის ან $1$-ის კოდირებისთვის სამ განსხვავებულ რიცხვს. შევნიშნოთ, რომ ბობს შეუძლია ორიგინალი რიცხვის გაშიფვრა ალისის მიერ გაგზავნილი სამი რიცხვიდან ნებისმიერი ორიდან.


\begin{table}[H]
\begin{tblr}{|Q[c,m]|Q[c,m]|Q[c,m]|}
\hline
\textbf{\makecell{შემსრულებელი}} & \textbf{\makecell{ფუნქციის გამოძახება}} & \textbf{დაბრუნებული მნიშვნელობა} \\
\hline
Alice & \texttt{setup(10)} & \texttt{2} \\
\hline
Bob & \texttt{setup(10)} & \texttt{2} \\
\hline
Alice & \texttt{encode(0)} & \texttt{\{5, 5, 5\}} \\
\hline
Alice & \texttt{encode(1)} & \texttt{\{8, 3, 7\}} \\
\hline
Alice & \texttt{encode(1)} & \texttt{\{0, 3, 1\}} \\
\hline
Alice & \texttt{encode(0)} & \texttt{\{7, 7, 7\}} \\
\hline
Alice & \texttt{encode(1)} & \texttt{\{6, 2, 0\}} \\
\hline
Bob & \texttt{decode(5, 5)} & \texttt{0} \\
\hline
Bob & \texttt{decode(8, 7)} & \texttt{1} \\
\hline
Bob & \texttt{decode(3, 0)} & \texttt{1} \\
\hline
Bob & \texttt{decode(7, 7)} & \texttt{0} \\
\hline
Bob & \texttt{decode(2, 0)} & \texttt{1} \\
\hline
\end{tblr}
\end{table}

\subsection{სანიმუშო გრეიდერი}

სანიმუშო გრეიდერისთვის, \texttt{encode}-სა და \texttt{decode}-ს ყველა გამოძახება თქვენი პროგრამის ერთსა და იმავე შესრულებაში განხორციელდება. გარდა ამისა, \texttt{setup} გამოიძახება მხოლოდ ერთხელ (განსხვავებით ორჯერ, ერთხელ შესრულებისას, როგორც შეფასების სისტემაშია).

შესატანი მონაცემი არის მხოლოდ ერთი მთელი რიცხვი -- $M$. შემდეგ ის დაბეჭდავს თქვენი \texttt{setup}-ის მიერ დაბრუნებულ $N$-ს. შემდეგ ის გამოიძახებს \texttt{encode} და \texttt{decode} ფუნქციებს ამ თანმიმდევრობით $T$ ჯერ, შემთხვევით გენერირებული რიცხვებით $0$-დან $N-1$-მდე და შემთხვევითად ირჩევს,  \texttt{encode}-ს დაბრუნებული სამი რიცხვიდან რომელი ორი უნდა გადასცეს \texttt{decode}-ს (და რა თანმიმდევრობით). ის დაბეჭდავს შეცდომის შეტყობინებას, თუ თქვენი ამოხსნა ვერ ამოხსნის მოცემულ ტესტს.

\end{document}
